\subsection{Data}
\label{sec:Data}
This study uses the CAMELS-US dataset (\href{https://doi.org/10.1016/0022-1694(89)90184-4}{https://doi.org/10.1016/0022-1694(89)90184-4}), which provides daily hydro-meteorological time series for minimally disturbed catchments across the contiguous United States\cite{CAMELS}. We select this dataset for three reasons: (1) catchments are weakly affected by direct human regulation, which helps identify more robust hydrological patterns; (2) long and standardized daily forcing--streamflow records support deep learning time-series training and cross-basin comparison\cite{hess-camels}; and (3) the data organization aligns with large-sample hydrology workflows, enabling reproducible evaluation\cite{nearing2021role}. To ensure temporal consistency across stations, we use a common analysis period from 1980-01-01 to 2014-12-31. Based on this period, a total of 521 stations are retained after data quality control and screening.The model input variables consist of meteorological forcing and antecedent hydrological information. Specifically, the meteorological forcing includes multiple daily atmospheric variables, while the antecedent hydrological information is represented by historical streamflow sequences to capture catchment state memory. Additional details can be found in the Supplementary Information \ref{sec:data}.

The prediction task is formulated as a sequence-to-one regression problem, where a sliding window of length $L=15$ days is used to predict next-day streamflow. All features are standardized using training-set statistics, and the data are split chronologically into training (70\%), validation (15\%), and test (15\%) subsets.

The model is implemented in PyTorch (Python 3.10) and trained on a CUDA-enabled Linux server. A three-stage pipeline (dense pre-training, weight sparsification, and mask learning) is optimized using AdamW with weight decay. Hyperparameters are tuned via Bayesian optimization\cite{BO_OPT} with validation NSE as the objective, and the selected configuration is consistently used across all stages and subsequent explainability analyses. Further details are provided in Supplementary Information \ref{sec:STEHhyper}.